% !PRINTED 280
\weberpage{280}
Wir können für $\varrho=\gamma^{12}$ eine beliebige primitive Wurzel von 11, z. B. 8, wählen. Dann wird $\varrho^{-1}=7$ und
\[
\varrho-\varrho^{-1}=1,\quad \varrho^2-\varrho^{-2}=4,
\]
also
\weq{13}{
\Theta = \begin{pmatrix} 3, & 0 \\ 0, & 4 \end{pmatrix}, \
\psi = \begin{pmatrix} 0, & 1 \\ -1, & 0 \end{pmatrix}, \
\chi = \begin{pmatrix} -3, & 1 \\ 1, & 3 \end{pmatrix}.
}
Damit ist also festgestellt, welche Arten von Theilern in der Gruppe $L_p$ vorkommen können, und es ist auch gezeigt, dass alle diese Theiler wirklich vorhanden sind. Die Frage, wie man für jeden Typus alle überhaupt möglichen Theiler erhält, haben wir hier bei Seite gelassen. Wir wollen darüber nur noch folgende Bemerkungen machen.

Wenn man eine gefundene Gruppe $G$ durch irgend eine Substitution der Gruppe $L_p$ transformirt, so erhält man einen conjugirten Theiler. Wenn man aber durch eine imaginäre Substitution der Gruppe $E_p$ transformirt, so kann es vorkommen, dass die transformirte Gruppe von $G$ trotzdem reell wird und nicht mit $G$ innerhalb $L_p$ conjugirt (wiewohl beide conjugirte Theiler von $E_p$ sind). So erhält man allgemein aus der Octaëder- und der Ikosaëdergruppe eine zweite nicht conjugirte, wenn man
mit $\begin{pmatrix} \varepsilon, & 0 \\ 0, & \varepsilon^{-1} \end{pmatrix}$ transformirt. Bei der Tetraëdergruppe giebt dies nur dann eine neue Gruppe, wenn $p\equiv\pm 1\ (\mathrm{mod}\ 8)$ ist. Die zweite Gruppe ergiebt sich nach der Formel:
\weq{14}{
\begin{pmatrix} \varepsilon, & 0 \\ 0, & \varepsilon^{-1} \end{pmatrix}
\begin{pmatrix} a, & b \\ c, & d \end{pmatrix}
\begin{pmatrix} \varepsilon^{-1}, & 0 \\ 0, & \varepsilon \end{pmatrix}
=
\begin{pmatrix} a, & bN \\ cN^{-1}, & d \end{pmatrix}
\quad (N=\varepsilon^2).
}
Um diese Verhältnisse an den Beispielen $p=5,7,11$ aufzuweisen, ist es zweckmässig, die oben gefundene Gruppe (6) für $p=7$ so zu transformiren, dass die Substitution $\chi$, die von der dritten Ordnung ist, die Normalform $S$ erhält. Man findet diese Transformation leicht, und erhält für diesen Fall:
\[
\begin{pmatrix} 1, & 1 \\ -2, & -1 \end{pmatrix}
\begin{pmatrix} 0, & 2 \\ 3, & 1 \end{pmatrix}
\begin{pmatrix} -1, & -1 \\ 2, & 1 \end{pmatrix}
=
\begin{pmatrix} 3, & 0 \\ 0, & 5 \end{pmatrix}
= \chi'
\]
\weq{15}{
\begin{pmatrix} 1, & 1 \\ -2, & -1 \end{pmatrix}
\begin{pmatrix} 2, & 1 \\ 3, & 2 \end{pmatrix}
\begin{pmatrix} -1, & -1 \\ 2, & 1 \end{pmatrix}
=
\begin{pmatrix} 1, & -2 \\ -1, & 3 \end{pmatrix}
= \Theta'
}
\[
\begin{pmatrix} 1, & 1 \\ -2, & -1 \end{pmatrix}
\begin{pmatrix} 2, & 2 \\ 1, & -2 \end{pmatrix}
\begin{pmatrix} -1, & -1 \\ 2, & 1 \end{pmatrix}
=
\begin{pmatrix} -3, & -3 \\ 1, & 3 \end{pmatrix}
= \psi',
\]
%CONT>>
